\section{Conclusion}
\label{sec:conclusion}

We rebuilt the project from its first report, retrained every model and evaluated on the full-resolution SID test set and on held-out ELD scenes. The answers to the three questions are more nuanced than the first version suggested.

\begin{enumerate}
\item \textbf{Bottleneck refinement with FiLM helps; teacher-latent distillation does not.} The Bridge U-Net without distillation is the best model (28.94\,dB, +0.32\,dB over our SID U-Net), and FiLM conditioning on the exposure ratio and noise level contributes 0.25\,dB. Distillation towards the encoder's own long-exposure features constrains the bridge and costs 0.58\,dB.
\item \textbf{A new camera needs very little data, and meta-learning helps only when adaptation is short.} One scene of an ELD camera adds about 2\,dB over the zero-shot model, and five scenes match full fine-tuning. Once support and query are separated by scene, MAML gives the best starting point and the best result after a few inner steps. With a longer adaptation the plain SID U-Net is as good or better.
\item \textbf{A real diffusion bridge is competitive but not better on fidelity metrics.} The feature-space variant of the first version is not a diffusion model and breaks under its own sampler. A pixel-space DDBM reaches 28.03\,dB in 5 steps, and more steps lower PSNR, as expected for posterior sampling. The released DDBM recipe trains far faster when its bridge noise is matched to the data range, but even then its few-step sampler leaves noise and it stays well below the simplified bridge.
\end{enumerate}

\paragraph{Limitations.} Every configuration was trained once. A second run of the SID U-Net differs by only \SIDSEEDGAP\,dB, but two runs of one model are a weak variance estimate, so the smaller gaps in Table~\ref{tab:sid} (for example FiLM, 0.25\,dB) need more seeds to be firm. The cloud machine running the first set of experiments was reclaimed before all files were copied. We kept every number reported here, but lost the checkpoints. The qualitative results therefore show only the two models retrained afterwards, and the $\sigma_{\max}=80$ DDBM has no test score. We measure fidelity only (PSNR, SSIM, MS-SSIM). A perceptual metric such as LPIPS or a user study would be needed to show the benefit of diffusion sampling.

\paragraph{Future work.} The strongest single result is that one or a few scenes are enough to adapt to a new sensor. This makes on-device adaptation feasible: a phone could capture a handful of short/long RAW pairs of a static scene and fine-tune the decoder and bridge only. Combining this with calibrated noise synthesis~\cite{wei2020eld}, and distilling the diffusion bridge into one or two steps, are the natural next steps.
